
\subsection{定理 7.3.1 的证明}
\label{subsec:7-3-2}
为实行上一小节末描述的计划，
还须重做第~\ref{sec:7-2} 节
论证中的一部分：再次扩大奇异链的可数集，
并克服第~\ref{subsec:7-2-3} 小节所述的困难。
不过，这两点都不会带来新问题。

先处理第一点。
我们将按 $g\in\mathbb Z_{\ge0}$ 归纳构造
一列集合 $\mathcal X_g(L)$。
（注意，起始的 $\mathcal X_0(L)$
已经通过类似的归纳在第~\ref{sec:7-2} 节构造
\translatornote{源书此处写“第 3.2 节”；
但旧单形集的逐代归纳在第 7.2 节，
紧接的“代”也明确回指第 7.2.4 小节，
故按定理 7.3.1 的陈述改引。}。
也就是说，
$\mathcal X_0(L)=\bigcup_g\mathcal X_{0,g}(L)$，
其中 $\mathcal X_{0,g}(L)$ 是
第~\ref{subsec:7-2-4} 小节意义下
第 $g$ 代的奇异单形集。
不过，在本节中，
把 $\mathcal X_0(L)$ 的全部元素都视为第零代。）

再处理第二点。
我们将证明 $\mathbf e$ 对 $A_{n,K}$ 结构
是同伦幺元。
这里
$(n,K)=(n(\lambda),K(\lambda))
\to(\infty,\infty)$，其中 $\lambda\to\infty$。
证明的其余部分便与第~\ref{subsec:7-2-8} 小节
的论证相似。
（严格说来，还需定义 $A_{n,K}$ 结构的
同伦幺元，并建立包含同伦幺元的
定理~\ref{thm:7-2-72} 的推广。
这项工作与第 4 章的讨论相似，
故原书从略。）

现在定义适当的模空间，
以开始证明定理~\ref{thm:7-3-1}。
取 $L$ 的光滑单纯剖分，置
\begin{equation}
[L]=\sum_c\pm\Delta_c^n,
\tag{7.3.6}\label{eq:7-3-6}
\end{equation}
其中 $n=\dim L$，右侧是该单纯剖分
中所有 $n$ 维单形的和。
(7.3.6) 中的符号 $\pm$ 如下确定：
先固定 $L$ 的三角剖分的顶点集次序，
据此把每个 $\Delta_c^n$ 识别为标准 $n$-单形，
从而赋予它定向。
当且仅当这一定向与 $L$ 在 $\Delta_c^n$
上诱导的定向相同，符号取 $+$。
可以假定 $L$ 的该单纯剖分的全部单形
都属于 $\mathcal X_0(L)$。

设
$\vec a=(a_1,\ldots,a_{|\vec a|})
\in\{1,\ldots,k\}^{|\vec a|}$，
其中 $a_1<\cdots<a_{|\vec a|}$。
令
\begin{equation}
\operatorname{forget}_{\vec a}\colon
\mathcal M^{\mathrm{main}}_{k+1}(\beta)
\longrightarrow
\mathcal M^{\mathrm{main}}_{k+1-|\vec a|}(\beta)
\tag{7.3.7.1}\label{eq:7-3-7-1}
\end{equation}
为遗忘第 $a_1,\ldots,a_{|\vec a|}$ 个
标记点的映射\liftIndex{I282}。
令
$\vec P=(P_1,\ldots,P_{k-|\vec a|})$，
其中 $P_i\in\mathcal X_{\mathfrak d(i)}(L)$。
在 $\vec P$ 中的第 $a_i$ 个位置
依次插入 $L$（$i=1,\ldots,|\vec a|$），
得到由 $k$ 条链组成的 $\vec P^+$。
例如，若 $\vec a=(1,3)$、$k=4$，
则 $\vec P^+=(L,P_1,L,P_2)$，
$\vec P=(P_1,P_2)$。
于是 (7.3.7.1) 诱导映射
\begin{equation}
\operatorname{forget}_{\vec a}\colon
\mathcal M^{\mathrm{main}}_{k+1}(\beta;\vec P^+)
\longrightarrow
\mathcal M^{\mathrm{main}}_{k+1-|\vec a|}(\beta;\vec P).
\tag{7.3.7.2}\label{eq:7-3-7-2}
\end{equation}

下面说明仓西结构的相容性。
取
\[
x\in\mathcal M^{\mathrm{main}}_{k+1-|\vec a|}
(\beta;\vec P),
\]
并设
$(V_x,E_x,\Gamma_x,\psi_x,s_x)$
为 $x$ 的仓西邻域。
也就是说，$V_x$ 是带有有限群
$\Gamma_x$ 作用的流形，
$E_x$ 是 $V_x$ 上的 $\Gamma_x$-等变丛，
$s_x$ 是 $E_x$ 的 $\Gamma_x$-等变截面，
而 $\psi_x$ 是从 $s_x^{-1}(0)/\Gamma_x$
到 $x$ 在
$\mathcal M^{\mathrm{main}}_{k+1-|\vec a|}
(\beta;\vec P)$ 中某邻域的同胚。
不妨设 $0\in V_x$、$s_x(0)=0$，
且 $\psi_x([0])=x$。

取
$x^+\in\mathcal M^{\mathrm{main}}_{k+1}
(\beta;\vec P^+)$，使
$\operatorname{forget}_{\vec a}(x^+)=x$，
并令
$(V_{x^+}^+,E_{x^+}^+,\Gamma_{x^+}^+,
\psi_{x^+}^+,s_{x^+}^+)$
为 $x^+$ 的仓西邻域。

\begin{lemma}{7.3.8}\label{lem:7-3-8}
可选取
$(V_{x^+}^+,E_{x^+}^+,\Gamma_{x^+}^+,
\psi_{x^+}^+,s_{x^+}^+)$，
使其满足以下性质：

\noindent\textup{(7.3.9.1)}\quad
$V_{x^+}^+\cong(W_x\times W'_{x^+})\times V_x$，
其中 $W_x\subset\mathbb R^{|\vec a|}$
是包含 $0$ 的开集，
$W'_{x^+}\subset\mathbb R^{c(x^+)}$
也是包含 $0$ 的开集；
$s_{x^+}^+(0)=0$，
$\psi_{x^+}^+([0])=x^+$。
这里 $c(x^+)$ 是依赖于 $x^+$ 的非负整数，
$\cong$ 表示分段光滑同构。

\noindent\textup{(7.3.9.2)}\quad
$\Gamma_{x^+}^+\cong\Gamma_x$。
它在第一个因子 $W_x\times W'_{x^+}$
上平凡作用；它在第二个因子 $V_x$
上的作用与 $\Gamma_x$ 的作用相同。

\noindent\textup{(7.3.9.3)}\quad
评价映射
$\operatorname{ev}_0^+\colon V_{x^+}^+\to L$
和 $\operatorname{ev}_0\colon V_x\to L$
满足
$\operatorname{ev}_0^+
=\operatorname{ev}_0\circ\pi$，
其中
$\pi\colon V_{x^+}^+
\cong(W_x\times W'_{x^+})\times V_x
\to V_x$
为到第二个因子的投影
\translatornote{源书此项将投影定义域写为
$W_x\times V_x$，漏了 (7.3.9.1) 中
的 $W'_{x^+}$；又把复合右侧写成
无下标的 $\operatorname{ev}\circ\pi$。
此处依完整积型邻域及已定义的
$\operatorname{ev}_0$ 补全。}。

\noindent\textup{(7.3.9.4)}\quad
在 $(W_x\times\{0\})\times V_x$ 上，
有
$E_{x^+}^+|_{(W_x\times\{0\})\times V_x}
\cong\pi^*E_x\oplus\mathbb R^{c(x^+)}$。

\noindent\textup{(7.3.9.5)}\quad
在子空间 $(W_x\times\{0\})\times V_x$ 上，
仓西映射 $s_{x^+}^+$ 的
$\pi^*E_x$ 分量是 $s_x$ 的拉回。

\noindent\textup{(7.3.9.6)}\quad
考察 $s_{x^+}^+$ 的
$\mathbb R^{c(x^+)}$ 分量。
在子空间 $(W_x\times\{0\})\times V_x$ 上，
它沿 $W'_{x^+}$ 方向的导数
等于自然嵌入
$T_0(W'_{x^+})\to
T_0(\mathbb R^{c(x^+)})
=\mathbb R^{c(x^+)}$。
\end{lemma}

\begin{proof}
以
$x^+=(\Sigma^+,\vec z^+,w^+)$
表示 $x^+$。
这里 $\Sigma^+$ 是亏格为零的
带边半稳定曲线，
$\vec z^+$ 是边界标记点，
而对给定的 $M$ 上殆复结构，
$w^+\colon(\Sigma^+,\partial\Sigma^+)
\to(M,L)$ 是伪全纯映射。
把遗忘 $|\vec a|$ 个标记点后所得的
三元组记作 $(\Sigma',\vec z',w')$；
这里 $\Sigma'$ 未必稳定。
此时 $x$ 由 $(\Sigma,\vec z,w)$ 表示，
其中 $\Sigma$ 是将 $\Sigma'$ 中
对应于 $(\Sigma',\vec z',w')$
不稳定分量的分量收缩后得到的。

若 $\Sigma'=\Sigma$，即
$(\Sigma',\vec z',w')$ 已经稳定，
则引理~\ref{lem:7-3-8}
实际上由构造便显然成立。
此时 $c(x^+)=0$。
（注意，构造仓西结构时，
我们并不破坏各种对称性。）

再考虑 $\Sigma'\ne\Sigma$ 的情形。
令 $\Sigma'=\bigcup_{i\in I}\Sigma_i$
为 $\Sigma'$ 的不可约分量分解，
并以 $I_0\subset I$ 表示其中
$\Sigma_i$ 不稳定的指标集。
依第 7.1 节构造仓西结构的一般方法，
$E_{x^+}^+$ 是某些有限维空间 $E_i$
的直和；$E_i$ 中的元素是
$(w^+)^*TM\otimes\Lambda^{(0,1)}\Sigma^+$
的光滑截面，
且对每个 $i\in I$，
各元素的支集都包含在
$\Sigma_i$ 的一个紧子集内。
在当前情形，可以选取仓西结构，
使经由
$(\Sigma',\vec z',w')$
向 $(\Sigma,\vec z,w)$ 的收缩，
有自然的同一
\[
\bigoplus_{i\in I\setminus I_0}E_i\cong E_x.
\]
现在置
$c(x^+)=\sum_{i\in I_0}\dim E_i$，
并定义 $W'_{x^+}$ 为
$\bigoplus_{i\in I_0}E_i$
中 $0$ 的一个小邻域。
回忆仓西映射 $s_{x^+}^+$、$s_x$
是先将 $\bar\partial$ 线性化，
再投影到各 $E_i$ 而得到的，
便容易核验 (7.3.9) 的各项性质。
例如，(7.3.9.6) 源于
所有 $i\in I_0$ 的分量 $\Sigma_i$
上的映射都是常值映射。
\end{proof}
